\subsubsection{Usando malloc}

Foi abordado no capítulo anterior a criação e manipulação de vetores estático, 
ou seja, que tinham um tamanho definido durante a sua declaração. No entanto, 
vamos supor que você tenha uma aplicação que contém um vetor, de qualquer tipo, 
com capacidade para guardar 12 itens. Porém, em algum momento da execução do 
programa surge a necessidade de se guardar um 13º item, dado a natureza dos 
vetores estáticos, seria necessário apagar algum dos itens existentes do vetor 
para guardar esse novo, caso contrário, seria impossível ter um 13º item dentro 
do vetor.

Esse problema, porém, pode ser resolvido facilmente com alocação de memória 
dinâmica. Através disso, é possível criar vetores dinâmicos que possuem um 
tamanho variável que pode ser modificado durante a execução da aplicação. 
Permitindo ao usuário uma maior flexibilidade e eficiência para armazenamento 
de certas informações.

Apesar dos pontos positivos, o uso de memória alocada também tem suas 
desvantagens. Como a alocação de memória não é feita de maneira automática pelo 
compilador durante o processo de compilação, e sim durante o processo de 
execução, é de responsabilidade do usuário cuidar daquele pedaço de memória e 
liberá-lo novamente para o sistema assim que possível. Processo esse que, caso 
não seja feito corretamente, pode causar problemas de vazamento de 
memória\footnote{Problemas de vazamento de memória são a causa de maior parte 
dos bugs e vulnerabilidades encontrados em softwares modernos, seções de 
memória alocadas e não liberadas podem ser a causa de mal funcionamento do 
sistema e até de ataques de agentes externas a sua aplicação/sistema.}.

Portanto, é necessário que ao utilizar memória alocada em sua aplicação o 
usuário lembre-se de liberar a memória após o seu uso. Ademais, vamos tentar 
exemplificar o processo de alocação de memória:

\begin{code}
    #include <stdio.h>
    #include <stdlib.h>

    int main(void)
    {
        // Declara um ponteiro para int
        int *numeros;
        
        int n; // Guarda uma quantidade de números a definir
        printf("Escolha a quantidade de números: ");
        scanf("%d", &n);
        
        // Reserva memória para 'n' inteiros
        numeros = malloc(n * sizeof(int));
        
        // Verifica se a memória foi alocada corretamente
        if (numeros == NULL) {
            printf("Alocação de memória falhou.\n");
            return 1;
        }
        
        printf("Digite %d números:\n", n);
        for (int i = 0; i < n; i++) {
            scanf("%d", &numeros[i]); // Recebe os números do usuário
        }
        
        // Calcula a soma
        int soma = 0;
        for (int i = 0; i < n; i++) {
            soma += numeros[i];
        }
        
        // Calcula a média
        float media = (float)soma / n;
        
        printf("Média: %.2f\n", media);
        
        // Libera a memória reservada pelo programa
        free(numeros);
        return 0;
    }
\end{code}

No código acima é pedido ao usuário para escolher uma quantidade \(N\) de 
números, que deverão ser digitados pelo usuário, de maneira que em seguida o 
programa calcule a média dos números digitados. Como o usuário pode escolher 
qualquer quantidade de números, não seria possível usar um vetor de tamanho 
estático.

(Para explicar o processo de alocação e liberação de memória vamos focar nas 
seguintes linhas de código)

\begin{code}
    numeros = malloc(n * sizeof(int));
    ...
    if (numeros == NULL) {
        printf("Alocação de memória falhou.\n");
        return 1;
    }
    ...
    free(numeros);
\end{code}

Sendo assim, declaramos um ponteiro \texttt{numeros} e atribuímos como valor o 
resultado da  função \texttt{malloc}, que aceita como parâmetro um número 
corresponde a quantidade de memória desejada a ser alocada e retorna o endereço 
do início da seção de memória alocada.

Após usar a função \texttt{malloc} para alocar a memória necessária, é 
necessário verificar se a memória foi alocada corretamente verificando se o 
ponteiro da seção de memória contém o valor \texttt{NULL} ou não. No exemplo 
acima verificamos se \texttt{numeros == NULL} e caso a verificação resulte em 
\texttt{verdadeiro} o programa imprime uma mensagem de erro ao usuário e é 
encerrado logo em seguida.

Após realizados todos os cálculos desejados usamos a função \texttt{free}, que 
aceita como parâmetro um ponteiro para a seção de memória alocada, para liberar 
a memória alocada devolta para o sistema.

\subsubsection{Malloc, Calloc e Realloc}

Aqui estão algumas funções similares a malloc que podem ser úteis:

\begin{itemize}
        \renewcommand{\labelitemi}{\tiny$\blacksquare$}
    \item \textbf{calloc:} Tem o mesmo funcionamento da função \texttt{malloc}, 
        porém, ao alocar a memória solicitada a função se certifica de 
        inicializar todo o bloco de memória alocado com o valor 0. Em alguns 
        casos seções de memória alocadas podem conter valores aleatórios que 
        foram deixados por outras por outras aplicações que usaram o bloco 
        anteriormente, então a depender do funcionamento da sua aplicação pode 
        ser necessário que se faça essa ``limpeza'' em blocos de memória 
        alocados. Segue abaixo um exemplo de como usar a função \texttt{calloc}:

        \begin{code}
    ...
    char *ponteiro = calloc(numero_de_itens, tamanho_de_cada_item);
    ...
        \end{code}
    \item \textbf{realloc:} A função realloc é a mais diferente dentre as três, 
        pois tem a função de re-alocar uma mesma seção de memória porém com um 
        tamanho diferente do anterior, geralmente essa função é utilizada 
        quando se utilizou uma seção de memória ao máximo e é necessário 
        expandi-la em algum outro momento. A função retorna um ponteiro para a 
        nova seção de memória com o tamanho modificado sem alterar o conteúdo 
        pré-existente daquela seção. Segue abaixo um exemplo de como usar a 
        função \texttt{realloc}:

        \begin{code}
    ...
    // Alocando memória para um vetor hipotético
    int *ponteiro = malloc(sizeof(int) * 10);
    // Modificando o tamanho da seção de memória alocada
    ponteiro = realloc(ponteiro, novo_tamanho_da_seção);
    ...
        \end{code}
\end{itemize}

Caso deseje aprender mais sobre o funcionamento das funções \texttt{malloc, 
realloc ou calloc}, é possível pesquisar utilizando o comando \texttt{`man 
nome\_da\_funcao'} em um terminal Unix/Linux, ou no seu navegador web, para 
acessar os manuais públicos relacionados a essas funções.
